% Kapitel 6 --- Ergebnisse
% Budget: 3 Seiten
% Regel: Jede Zahl kommt aus einem Skriptlauf, nicht aus dem Kopf.

\section{Ergebnisse}
\label{sec:ergebnisse}

\subsection{Abweichungsraten je Sachverhaltstyp}

\todo{Tabelle mit Typ, Anzahl Sachverhalte, Abweichungen, Anteil, 95-Prozent-Intervall.
Ein Satz je Zeile, der die Zahl einordnet. Die Tabelle ist das Ergebnis, nicht die
Vorbereitung darauf.}

\subsection{Verfahrensvergleich}

\todo{Nur Sprachmodell, nur Regeln, Hybrid: Precision, Recall, F1 gegen denselben
Goldstandard. Aussage: Was bringt die Normalisierung wirklich?}

\subsection{Fehleranalyse}

\todo{20 bis 30 Fehlfälle, kategorisiert und gezählt. Diese Tabelle zeigt, dass das
eigene Verfahren verstanden wird --- sie ist oft der stärkste Abschnitt einer Arbeit.}

\subsection{Beispiele}

\todo{Zwei Fälle ausgeschrieben: ein echter Widerspruch und ein Versionsartefakt, jeweils
mit Fundstellen und Typzuordnung. Nicht mehr als zwei.}
